% from ega4-11-fr.tex
\oldpage[IV-3]{158}

aussi bijectif (\emph{loc. cit.}, n$^\circ$~2, cor. de la prop.~6) ce qui est
absurde. Autrement dit, si l'on pose $Y=\operatorname{Spec}(A)$,
$X=\operatorname{Spec}(B)$, $\mathcal O_X$ n'est pas $Y$-plat en l'unique point
$x$ de $X$. Mais soit $A_1=B$, et $Y_1=\operatorname{Spec}(A_1)$, et considérons
l'homomorphisme canonique $A\to A_1$, qui est local et injectif puisque
$\mathfrak J\cap A=0$~; si
$B_1=B\otimes_A A_1=B\otimes_A B$, nous allons voir qu'il y a un point de
$X_1=X\times_Y Y_1=\operatorname{Spec}(B_1)$ où $\mathcal O_{X_1}$ est
$Y_1$-plat. Pour cela, remarquons que l'on a
\[
B\otimes_A B
=\frac{A'\otimes_A A'}
{\operatorname{Im}(\mathfrak J\otimes_A A')
+\operatorname{Im}(A'\otimes_A\mathfrak J)}.
\]
Or la structure de $C'=A'\otimes_A A'$ s'obtient aisément~; on considère la
$A$-algèbre produit $C=\prod_{\sigma\in\Gamma}A'_{\sigma}$, où tous les
$A'_{\sigma}$ sont égaux à $A'$, et l'application canonique
$\varphi:A'\otimes_A A'\to C$ telle que
$\varphi(x\otimes y)=(\sigma(x)y)_{\sigma\in\Gamma}$ (le groupe $\Gamma$
opérant canoniquement sur $A'$)~; en passant aux quotients, on en déduit un
homomorphisme $C'/\mathfrak m C'\to C/\mathfrak m C$ qui n'est autre que
l'homomorphisme canonique
$k'\otimes_k k'\to\prod_{\sigma}k'_{\sigma}$ (avec $k'_{\sigma}=k'$ pour tout
$\sigma\in\Gamma$)~; on sait que ce dernier est bijectif (Bourbaki,
\emph{Alg.}, chap.~VIII, \S~8, prop.~4), donc il en est de même de $\varphi$,
puisque $C'$ et $C$ sont des $A$-modules libres (Bourbaki,
\emph{Alg. comm.}, chap.~II, \S~3, n$^\circ$~2, cor. de la prop.~6). De ce qui
précède, il résulte que $B\otimes_A B$ est une $A$-algèbre semi-locale, composée
directe des $A$-algèbres locales $A'/(\sigma(\mathfrak J)+\mathfrak J)$. Celle
de ces algèbres correspondant à l'identité de $\Gamma$ est isomorphe à
$A'/\mathfrak J$, donc est un $A_1$-module plat, mais comme $\gamma\notin k$, il
y a au moins un $\sigma\in\Gamma$ tel que
$\sigma(\mathfrak J)\ne\mathfrak J$~; alors
$\sigma(\mathfrak J)+\mathfrak J=\mathfrak m'$, et $A'/\mathfrak m'$ n'est pas
un $(A'/\mathfrak J)$-module plat, puisque c'est le quotient de
$A'/\mathfrak J$ par un idéal non nul.
\end{env}

\begin{env}[11.7.2]
\label{IV.11.7.2-fr}
Nous allons maintenant montrer que le résultat de (11.5.4) perd sa validité
lorsqu'on ne suppose plus que $f$ soit un morphisme \emph{propre} (et
\emph{a fortiori} (11.5.3) cesse d'être exact lorsqu'on ne suppose plus $g$
propre). Soient $k$ un corps, $A_0$ l'anneau de polynômes $k[S,T]$, $A$
l'anneau quotient $A_0/A_0ST$~; $Y=\operatorname{Spec}(A)$ est donc la courbe
réductible formée des deux «~axes de coordonnées~» dans le plan affine
$V_k^2=\operatorname{Spec}(A_0)$. L'anneau $A$ admet deux idéaux premiers
minimaux $\mathfrak p_1=A_0S/A_0ST$, $\mathfrak p_2=A_0T/A_0ST$, et comme $A$
est réduit, il se plonge canoniquement dans $B=B_1\oplus B_2$, où
$B_1=A/\mathfrak p_1$, $B_2=A/\mathfrak p_2$~; d'ailleurs $B_1$ s'identifie à
$k[T]$ et $B_2$ à $k[S]$, donc ce sont des anneaux intègres intégralement clos
et par suite $Z=\operatorname{Spec}(B)$ n'est autre que le normalisé du
préschéma $Y$ par rapport à $R(Y)$ (\textbf{II}, 6.3.8), somme des deux schémas
$Z_1=\operatorname{Spec}(B_1)$, $Z_2=\operatorname{Spec}(B_2)$. Désignons par
$y$ le «~point double~» de $Y$, correspondant à l'idéal maximal
$\mathfrak p_1+\mathfrak p_2=\mathfrak m$ de $A$, par $z_1$ et $z_2$ les
points de $Z$ qui se projettent en $y$, correspondant respectivement aux idéaux
maximaux $\mathfrak m_1=(T)$ et $\mathfrak m_2=(S)$ de $B_1$ et $B_2$. Nous
désignerons par $X$ le sous-préschéma de $Z$ induit sur le complémentaire de
$z_2$ dans $Z$~; on a donc
$X=\operatorname{Spec}(B_1\oplus(B_2)_S)$~; il est immédiat que
l'homomorphisme $A\to A'=B_1\oplus(B_2)_S$ est injectif, mais le morphisme
correspondant $f:X\to Y$ n'est pas fermé (car $f(Z_2-\{z_2\})$ n'est pas fermé
dans $Y$, bien que $Z_2-\{z_2\}$ soit fermé dans $X$)~; \emph{a fortiori} il
n'est pas propre. Nous allons voir maintenant que $f$ n'est pas plat au point
$z_1$~; il suffira de montrer $(\mathbf 0_{\mathrm I},6.6.2)$ que
$\mathcal O_{z_1}$ n'est pas un $\mathcal O_y$-module fidèlement plat, et pour
cela il suffit de voir que l'homomorphisme canonique
$\mathcal O_y\to\mathcal O_{z_1}$ n'est pas injectif~; mais cela est immédiat
car $\mathcal O_{z_1}$ est un anneau intègre, tandis que $\mathcal O_y$ admet
des diviseurs de zéro. Cependant, la première projection
$p:X\times_Y X\to X$ est un isomorphisme~: en effet, on a
\[
B_1\otimes_A B_1=(A/\mathfrak p_1)\otimes_A(A/\mathfrak p_1)=A/\mathfrak p_1,
\]
$(B_2)_S\otimes_A(B_2)_S=(B_2\otimes_A B_2)_S=(B_2)_S$ pour la même raison,
et enfin $B_1\otimes_A(B_2)_S=0$, car l'image canonique de $S$ dans $B_1$ est
nulle.
\end{env}

\begin{env}[11.7.3]
\label{IV.11.7.3-fr}
L'exemple précédent peut se généraliser~: on considère sur un corps $k$ une
courbe algébrique réduite $Y$ admettant un seul «~point double ordinaire~» $y$
(notion qui sera définie plus tard de façon générale), et sa normalisée $Z$, de
sorte que le morphisme $g:Z\to Y$ est fini, que la restriction de $g$ à
$Z-g^{-1}(y)$ est un isomorphisme sur $Y-\{y\}$, et que $g^{-1}(y)$ se réduit à
deux points «~simples~» $z_1,z_2$~; en outre le préschéma $g^{-1}(y)$ est somme
des deux préschémas $\operatorname{Spec}(k(z_1))$,
$\operatorname{Spec}(k(z_2))$, canoniquement isomorphes à
$\operatorname{Spec}(k(y))$. Soit $X$ le sous-préschéma de $Z$ induit sur
l'ouvert $Z-\{z_2\}$~; le morphisme $f:X\to Y$, restriction de $g$ à $X$, n'est
pas propre, sans quoi (\textbf{II}, 5.4.3) il en serait de même de l'injection
canonique $j:X\to Z$, qui n'est pas fermée. Le morphisme $f$ est radiciel, car
pour tout $y'\in Y$, la fibre $f^{-1}(y')$ ne comprend qu'un seul point $x'$,
$k(y')\to k(x')$ étant un isomorphisme~; on en conclut d'abord que le morphisme
diagonal $\Delta_f:X\to X\times_Y X$ est bijectif (1.7.7.1) et d'autre part,
comme $f$ est non ramifié (17.4.2, d$'$), $\Delta_f$ est une immersion ouverte
(17.4.2, b$'$)~; par suite $\Delta_f$ est un isomorphisme, et la première
projection $p:X\times_Y X\to X$ l'isomorphisme réciproque. Cependant $f$ n'est
pas plat au point $z_1$~; sinon $\mathcal O_{z_1}$ serait un
$\mathcal O_y$-module fidèlement plat $(\mathbf 0_{\mathrm I},6.6.2)$, et comme
$\mathcal O_y$ contient deux idéaux premiers minimaux $\mathfrak p_1$,
$\mathfrak p_2$ distincts (correspondant aux deux «~branches~» de $Y$ au point
$y$) il existerait dans $\mathcal O_{z_1}$ deux idéaux premiers dont les images
réciproques par $u:\mathcal O_y\to\mathcal O_{z_1}$ seraient
$\mathfrak p_1$ et $\mathfrak p_2$ $(\mathbf 0_{\mathrm I},6.5.1)$~; mais cela
est absurde, car $\mathcal O_{z_1}$ n'a que deux idéaux premiers distincts, $0$
et l'idéal maximal $\mathfrak m_1$, et $u^{-1}(\mathfrak m_1)$ est l'idéal
maximal $\mathfrak m$ de $\mathcal O_y$.
\end{env}

\begin{env}[11.7.4]
\label{IV.11.7.4-fr}
On notera que dans l'exemple précédent l'homomorphisme $u$ est \emph{injectif}
lorsque $Y$ est \emph{irréductible} (on peut par exemple prendre
$Y=\operatorname{Spec}(k[S,T]/(S(S^2+T^2)-(S^2-T^2)))$, «~cubique à point
double~»)~; on peut en effet alors (en remplaçant $Y$ par un voisinage affine de
$y$) supposer que $Y=\operatorname{Spec}(A)$, où $A$ est intègre, d'où
$Z=\operatorname{Spec}(B)$, où $B$ est la clôture intégrale de $A$~; comme $B$,
$\mathcal O_y$ et $\mathcal O_{z_1}$ s'identifient alors à des sous-anneaux du
corps
% IV3_B23_P158_END

